\documentclass[letter,12pt]{article}

% --- Packages (matching original JPE/JEP style) ---
\usepackage{amsfonts, amssymb, amsmath, bbold}
\usepackage{booktabs, dcolumn, multirow, array, tabularx, tabu, threeparttable}
\usepackage{float, pdflscape, geometry, setspace}
\usepackage{mathpazo}
\usepackage[dvipsnames]{xcolor}
\usepackage{caption, subcaption}
\usepackage{graphicx}
\usepackage{hyperref}
\usepackage{natbib}

\geometry{left=1.25in, right=1.25in, top=1.25in, bottom=1.25in}
\doublespacing

\renewcommand{\arraystretch}{1.3}
\def\sym#1{\ifmmode^{#1}\else\(^{#1}\)\fi}

% Extra row separation for tabu
\extrarowsep=0.1pt

% Significance star macro
\newcommand{\starnote}{
  \noindent\emph{Notes}: Asterisks indicate statistical significance (2-tailed):
  * $P < 0.10$; ** $P < 0.05$; *** $P < 0.01$.
  Robust standard errors clustered at the community level in parentheses.
}

\hypersetup{colorlinks=true, linkcolor=black, citecolor=black, urlcolor=black}

\pagenumbering{arabic}

% ------------------------------------------------------------------
\begin{document}

% ------------------------------------------------------------------
%  TITLE
% ------------------------------------------------------------------
\begin{titlepage}
\begin{center}
\vspace*{2cm}
{\Large \textbf{Do Expectations Shape Wellbeing Gains from Electrification?\\
Willingness-to-Pay Heterogeneity and Life Satisfaction\\
in Rural Kenya}}\\[2em]
{\large Extension of Lee, Miguel \& Wolfram (JPE 2019; JEP 2019)}\\[2em]
{\normalsize \today}\\[6em]

\begin{abstract}
\noindent
This paper asks whether baseline willingness-to-pay (WTP) for a grid electricity
connection predicts the life satisfaction gains from electrification. Using data
from a randomised experiment in rural Kenya, we merge stated WTP elicited at
baseline (2014) with subjective wellbeing outcomes measured at two follow-up
rounds (2016, 2017). We estimate heterogeneous treatment-on-treated (TOT)
effects on a life satisfaction index and a mental health symptoms index,
using the three experimental subsidy arms as instruments and WTP as the
interaction variable, mirroring the specification of Lee, Miguel \& Wolfram
(2019). We also examine whether, conditional on connecting, realised life
satisfaction scales with stated WTP. Our findings speak to the role of
prior expectations in shaping welfare responses to household electrification.
\end{abstract}
\end{center}
\end{titlepage}

% ------------------------------------------------------------------
%  1. INTRODUCTION
% ------------------------------------------------------------------
\section{Introduction}

A large literature on rural electrification has documented modest average
impacts on household welfare. Lee, Miguel \& Wolfram (2019) find small or
null average treatment effects on life satisfaction in a randomised experiment
in Kenya, even as connection rates respond strongly to price subsidies. One
explanation is that mean effects mask substantial heterogeneity: households
with higher prior expectations of benefit may experience larger wellbeing
gains, while those who connected despite low willingness-to-pay may see
little change.

We test this mechanism directly. At baseline, each sampled household was
presented with a randomly assigned hypothetical connection price and asked
whether they would accept. This contingent valuation elicitation provides a
measure of stated WTP that varies across households. We link this to
subjective wellbeing outcomes at two follow-up rounds to ask: do households
with higher baseline WTP experience larger life satisfaction gains from
electrification?

We estimate two complementary analyses:

\begin{enumerate}
  \item \textbf{Heterogeneous TOT by baseline WTP.} We instrument household
    grid connection with the three experimental treatment arms and their
    interactions with a high-WTP indicator, following the specification of
    Lee, Miguel \& Wolfram (2019, Appendix Tables A3--A4).

  \item \textbf{WTP as a predictor of life satisfaction conditional on
    connecting.} Among connected households, we regress life satisfaction
    on stated WTP (in USD), controlling for the same household and community
    covariates.
\end{enumerate}

% ------------------------------------------------------------------
%  2. DATA AND IDENTIFICATION
% ------------------------------------------------------------------
\section{Data and Identification}

\subsection{Data}

We use three datasets from Lee, Miguel \& Wolfram (2019):

\begin{itemize}
  \item \texttt{demand.dta} --- 2014 baseline household survey with WTP
    elicitation and treatment assignment.
  \item \texttt{impacts\_1.dta} --- 2016 Round 1 (R1) follow-up survey with
    life satisfaction and mental health outcomes.
  \item \texttt{impacts\_2.dta} --- 2017 Round 2 (R2) follow-up survey with
    the same outcomes.
\end{itemize}

The three datasets share a unique household identifier (\texttt{reppid}),
enabling a 1:1 merge of baseline WTP with follow-up wellbeing outcomes.
We exclude households already connected at baseline
(\texttt{base\_connected}~$=$~1).

\subsection{WTP Measure}

Each household was randomly assigned a contingent valuation price
(\texttt{WTP\_amt} $\in \{0, 10\text{k}, 15\text{k}, 20\text{k}, 25\text{k},
35\text{k}, 75\text{k}\}$ KES) and asked whether they would accept. We
define:

\begin{equation}
  \texttt{wtp\_stated}_i =
    \texttt{WTP\_amt}_i \times \mathbf{1}[\texttt{WTP\_r1}_i = 1]
\end{equation}

and \texttt{wtp\_high}$_i = 1$ if \texttt{wtp\_stated}$_i$ exceeds the
sample median. Because \texttt{WTP\_amt} is randomly assigned, stated WTP
has an exogenous component that is uncorrelated with unobservables.

\subsection{Outcomes}

Primary outcome: \texttt{life\_index} (normalised life satisfaction index).
Secondary outcome: \texttt{symptoms\_index} (mental health symptoms index,
higher values indicate worse symptoms). Both are available in R1 and R2.

\subsection{Identification}

\paragraph{Analysis 1 --- Heterogeneous TOT.}
We estimate:
\begin{equation}
  Y_{ic} = \alpha + \beta_1 \widehat{D}_{ic} + \beta_2 (\widehat{D}_{ic}
    \times \text{WTP-High}_{i}) + \gamma \, \text{WTP-High}_{i}
    + \delta' \mathbf{X}_{ic} + \varepsilon_{ic}
\end{equation}
where $D_{ic}$ is an indicator for grid connection, instrumented by
\texttt{treat\_low}, \texttt{treat\_med}, \texttt{treat\_full} and their
interactions with \texttt{wtp\_high}. Standard errors are clustered at the
community level. $\beta_1$ is the TOT for low-WTP households;
$\beta_1 + \beta_2$ is the TOT for high-WTP households.

\paragraph{Analysis 2 --- Conditional OLS.}
Among connected households at R1 follow-up we estimate:
\begin{equation}
  Y_{ic} = \alpha + \beta \, \text{WTP-Stated}_{i} (\text{USD})
    + \delta' \mathbf{X}_{ic} + \varepsilon_{ic}
\end{equation}
This is descriptive: $\beta$ captures whether ex-ante WTP predicts
ex-post life satisfaction among those who connected, consistent with an
expectations channel but not cleanly causal due to selection into connection.

% ------------------------------------------------------------------
%  3. FIGURES
% ------------------------------------------------------------------
\section{Demand Curves and WTP Heterogeneity at Baseline}

Figure~\ref{fig:demand_income} reproduces the experimental demand curves
from Lee, Miguel \& Wolfram (2019), splitting the sample by baseline income
quartile (panel A) and housing quality (panel B). High-income and high-quality
households reveal higher WTP through greater take-up at any given price ---
the same households we expect to benefit more from connection if expectations
drive wellbeing gains.

\begin{figure}[H]
  \centering
  \begin{subfigure}[t]{0.47\textwidth}
    \centering
    \includegraphics[width=\linewidth]{../../data/jpe/figures/figure2b.png}
    \caption{Demand curve by baseline income (bottom vs.\ top earnings quartile)}
    \label{fig:demand_income_a}
  \end{subfigure}
  \hfill
  \begin{subfigure}[t]{0.47\textwidth}
    \centering
    \includegraphics[width=\linewidth]{../../data/jpe/figures/figure2c.png}
    \caption{Demand curve by housing quality (low vs.\ high wall quality)}
    \label{fig:demand_income_b}
  \end{subfigure}
  \caption{\textbf{Demand for grid connection by baseline household
    characteristics.}
    Solid lines show take-up rates for the high group; dashed lines for
    the low group. Price on the horizontal axis is in USD. Each point is
    the share of households that took up the connection at the assigned
    experimental price. Source: Lee, Miguel \& Wolfram (2019, Figure 2).}
  \label{fig:demand_income}
\end{figure}

% ------------------------------------------------------------------
%  4. RESULTS
% ------------------------------------------------------------------
\section{Results}

\subsection{Heterogeneous Treatment Effects on Wellbeing by Baseline WTP}

Tables~\ref{tab:wtp_r1} and~\ref{tab:wtp_r2} report IV estimates of the
treatment effect on life satisfaction and mental health, split by high vs.\
low baseline WTP, for R1 (2016) and R2 (2017) respectively.

The coefficient on \texttt{connected} gives the TOT for low-WTP households.
The coefficient on \texttt{connected}~$\times$~\texttt{wtp\_high} gives the
\emph{differential} TOT for high-WTP households relative to low-WTP
households. A positive and significant differential would indicate that
high-WTP households benefit more from electrification in terms of life
satisfaction, consistent with an expectations channel.

\clearpage
\newgeometry{left=1in, right=1in, top=1in, bottom=1in}

\begin{table}[H]
  \centering
  \caption{\textbf{Heterogeneous TOT on Wellbeing by Baseline WTP --- R1 (2016)}}
  \label{tab:wtp_r1}
  \vspace{0.5em}
  \input{../tables/tab_wtp_r1.tex}
  \vspace{0.5em}
  \begin{minipage}{\linewidth}
    \small
    \emph{Notes}: IV-2SLS estimates. Dependent variables are
    \texttt{life\_index} (column 1) and \texttt{symptoms\_index} (column 2).
    \texttt{Connected} is instrumented by \texttt{treat\_low},
    \texttt{treat\_med}, \texttt{treat\_full} and their interactions with
    \texttt{wtp\_high}. \texttt{wtp\_high} is an indicator for stated WTP
    above the sample median. Controls include community-level covariates
    (county, market proximity, transformer timing, baseline connection rate,
    population) and household covariates (gender, age, housing quality,
    baseline energy spending). Sample restricted to adult household members
    (\texttt{child}~$=$~0) not connected at baseline.
    \starnote
  \end{minipage}
\end{table}

\clearpage

\begin{table}[H]
  \centering
  \caption{\textbf{Heterogeneous TOT on Wellbeing by Baseline WTP --- R2 (2017)}}
  \label{tab:wtp_r2}
  \vspace{0.5em}
  \input{../tables/tab_wtp_r2.tex}
  \vspace{0.5em}
  \begin{minipage}{\linewidth}
    \small
    \emph{Notes}: As Table~\ref{tab:wtp_r1} but using R2 (2017) follow-up data.
    \starnote
  \end{minipage}
\end{table}

\restoregeometry

\subsection{WTP as a Predictor of Life Satisfaction Conditional on Connecting}

Table~\ref{tab:wtp_connected} presents OLS estimates of the relationship
between stated WTP (in USD, 2014 prices) and life satisfaction, restricted
to households connected at the R1 follow-up.

\clearpage
\newgeometry{left=1in, right=1in, top=1in, bottom=1in}

\begin{table}[H]
  \centering
  \caption{\textbf{WTP as Predictor of Wellbeing Conditional on Connection --- R1 (2016)}}
  \label{tab:wtp_connected}
  \vspace{0.5em}
  \input{../tables/tab_wtp_connected.tex}
  \vspace{0.5em}
  \begin{minipage}{\linewidth}
    \small
    \emph{Notes}: OLS estimates. Sample restricted to households connected to
    the grid at R1 (\texttt{connected}~$=$~1). \texttt{wtp\_stated\_usd} is
    the stated WTP from the 2014 contingent valuation, converted to USD at
    the 2014 average exchange rate (87.94 KES/USD). Controls are as in
    Table~\ref{tab:wtp_r1}. This regression is descriptive: $\beta$ captures
    whether ex-ante WTP predicts ex-post life satisfaction among connectors,
    but is not a clean causal estimate due to selection into connection.
    \starnote
  \end{minipage}
\end{table}

\restoregeometry

\subsection{Robustness}

Table~\ref{tab:robustness} presents three robustness checks using R1 data:
(i) replacing \texttt{wtp\_high} with a binary indicator for any positive WTP
under the 6-week financing window (\texttt{wtp\_r2\_any}), (ii) using
\texttt{wtp\_any} (willing to pay anything at all), and (iii) including both
\texttt{wtp\_high} and the SES composite index (\texttt{highses}) from
Lee, Miguel \& Wolfram (2019) to test whether WTP heterogeneity simply
reflects socioeconomic status.

\clearpage
\newgeometry{left=1in, right=1in, top=1in, bottom=1in}

\begin{table}[H]
  \centering
  \caption{\textbf{Robustness: Alternative WTP Measures and SES Control --- R1 (2016)}}
  \label{tab:robustness}
  \vspace{0.5em}
  \input{../tables/tab_wtp_robustness.tex}
  \vspace{0.5em}
  \begin{minipage}{\linewidth}
    \small
    \emph{Notes}: Dependent variable is \texttt{life\_index} throughout.
    Column (1) uses \texttt{wtp\_r2\_any} (binary WTP under 6-week financing,
    \texttt{WTP\_r2}). Column (2) uses \texttt{wtp\_any} (any positive stated WTP
    from \texttt{WTP\_r1}). Column (3) includes both \texttt{wtp\_high} and
    \texttt{highses} (SES composite index above 75th percentile, constructed
    as in Lee, Miguel \& Wolfram 2019) to test whether WTP captures SES
    rather than expectations. All specifications are IV-2SLS with community-
    clustered standard errors. Sample as in Table~\ref{tab:wtp_r1}.
    \starnote
  \end{minipage}
\end{table}

\restoregeometry

% ------------------------------------------------------------------
%  5. CONCLUSION
% ------------------------------------------------------------------
\section{Conclusion}

We examine whether willingness-to-pay for grid electricity at baseline
predicts life satisfaction gains at follow-up in a randomised experiment in
rural Kenya. By substituting the baseline WTP indicator for the SES composite
used in Lee, Miguel \& Wolfram (2019), we test an expectations channel:
households that valued the connection more ex-ante may also benefit more
ex-post.

Our findings [to be reported upon running the Stata code] shed light on
whether the well-documented null average effects on wellbeing conceal
heterogeneity along the dimension of prior expectations.

% ------------------------------------------------------------------
%  REFERENCES
% ------------------------------------------------------------------
\clearpage
\begin{singlespace}
\bibliographystyle{aer}
\begin{thebibliography}{9}

\bibitem{lee2019jpe}
Lee, Kenneth, Edward Miguel, and Catherine Wolfram. 2019.
``Experimental Evidence on the Economics of Rural Electrification.''
\textit{Journal of Political Economy} 128 (4): 1523--1565.

\bibitem{lee2019jep}
Lee, Kenneth, Edward Miguel, and Catherine Wolfram. 2020.
``Does Household Electrification Supercharge Economic Development?''
\textit{Journal of Economic Perspectives} 34 (1): 122--144.

\end{thebibliography}
\end{singlespace}

\end{document}
